% !TEX root = ../main.tex
% !TEX program = lualatex
\documentclass[../main.tex]{subfiles}
\IfSubfilesClassLoaded{\externaldocument{\subfix{../build/main}}}{}
% ====================
% File: 17K_jcids.tex
% ====================
\begin{document}
\IfSubfilesClassLoaded{\chapter{EDO Study Guide}}{}
% === JCIDS (EDO 3.3.3) ===
\section{Joint Capabilities Integration \& Development System}

\subsection{Summary}
\begin{description}
  \item[\textbf{Decision support triad.}] \ac{jcids} (requirements), \ac{ppbe} (funding), and the \ac{das} (materiel) operate continuously and must stay aligned to enable timely decisions~\autocite{edo-3-3-3-jcids-2025}.
  \item[\textbf{Current JFRP note.}] Current Joint Staff guidance cancels the 2021 \ac{jcids} manual and rescinds/replaces the \ac{jcids} process with the Joint Force Requirements Process (JFRP); this section preserves \ac{jcids} terminology for legacy EDO lecture context and calls out current JFRP governance where it changes the answer~\autocite{CJCSM5123-01}.
  \item[\textbf{Purpose.}] \ac{jcids} ensures the joint warfighter receives required capabilities, institutionalizes joint needs, and initiates the acquisition management process~\autocite{edo-3-3-3-jcids-2025}.
  \item[\textbf{CBA foundation.}] \acp{cba} identify capability gaps and recommend materiel or non-materiel solutions based on operational risk and feasibility~\autocite{edo-3-3-3-jcids-2025}.
  \item[\textbf{Capability documents.}] \ac{icd} and \ac{cdd} (with \ac{cdd} updates) align to acquisition phases and provide validated requirements for program execution~\autocite{edo-3-3-3-jcids-2025}.
  \item[\textbf{Key players.}] \ac{cjcs}, \ac{jroc}, \ac{jcb}, \acp{fcb}, the Joint Staff Gatekeeper (J-8), and Joint Staff J-6 drive validation, prioritization, and interoperability decisions~\autocite{edo-3-3-3-jcids-2025}.
\end{description}

\subsection{Decision Support Systems}
\begin{description}
  \item[\textbf{Requirements.}] \ac{jcids} captures warfighter capability needs and requirements~\autocite{edo-3-3-3-jcids-2025}.
  \item[\textbf{Funding.}] \ac{ppbe} provides the resource decisions to execute validated priorities~\autocite{edo-3-3-3-jcids-2025}.
  \item[\textbf{Materiel.}] The \ac{das} executes acquisition strategies that deliver validated capabilities~\autocite{edo-3-3-3-jcids-2025}.
  \item[\textbf{Synchronization.}] These systems must interface continuously to keep requirements, funding, and execution aligned~\autocite{edo-3-3-3-jcids-2025}.
\end{description}

\begin{figure}[htb]
  \centering
  \includegraphics[width=0.8\linewidth]{JCIDStoDAS.png}
  \caption[JCIDS to DAS Relationship]{Relationship between \ac{jcids} and the \ac{das}. \srcCite{edo-3-3-3-jcids-2025}.}
  \label{fig:jcids-to-das}
\end{figure}

\subsection{Capabilities-Based Assessment (CBA)}
\begin{description}
  \item[\textbf{Objective.}] A \ac{cba} analyzes mission needs and capability gaps, assesses risk, and recommends materiel and non-materiel solutions~\autocite{edo-3-3-3-jcids-2025}.
  \item[\textbf{Initiation.}] A \ac{cba} can start from any \ac{dod} organization, a joint concept, a \ac{conops}, or an operational need~\autocite{edo-3-3-3-jcids-2025}.
  \item[\textbf{Outputs.}] The \ac{cba} identifies mission and capability needs, gaps and risks, validates non-materiel feasibility, and recommends a path forward~\autocite{edo-3-3-3-jcids-2025}.
\end{description}

\subsection{CBA Phases}
\begin{description}
  \item[\textbf{Study definition.}] What does the mission require~\autocite{edo-3-3-3-jcids-2025}?
  \item[\textbf{Needs assessment.}] How well do current capabilities meet those needs; where are the gaps~\autocite{edo-3-3-3-jcids-2025}?
  \item[\textbf{Solutions recommendation.}] What materiel or non-materiel options should close the gaps~\autocite{edo-3-3-3-jcids-2025}?
\end{description}

\subsection{CBA Recommended Solutions}
\begin{description}
  \item[\textbf{Materiel.}] New or modified systems acquired through the \ac{das}; guided by validated \ac{jcids} documents~\autocite{edo-3-3-3-jcids-2025}.
  \item[\textbf{Non-materiel.}] Changes in doctrine, organization, training, leadership, personnel, facilities, or policy; often faster and less costly than new systems~\autocite{edo-3-3-3-jcids-2025}.
  \item[\textbf{Document path.}] Non-materiel solutions lead to a \ac{dcr}; materiel solutions lead to an \ac{icd}~\autocite{edo-3-3-3-jcids-2025}.
\end{description}

\subsection{DOTMLPF-P Elements}
\begin{description}
  \item[\textbf{Doctrine.}] Fundamental guiding principles~\autocite{edo-3-3-3-jcids-2025}.
  \item[\textbf{Organization.}] How the \ac{dod} organizes to fight~\autocite{edo-3-3-3-jcids-2025}.
  \item[\textbf{Training.}] How the force prepares to fight tactically~\autocite{edo-3-3-3-jcids-2025}.
  \item[\textbf{Materiel.}] Existing equipment and systems~\autocite{edo-3-3-3-jcids-2025}.
  \item[\textbf{Leadership and education.}] Leader development requirements~\autocite{edo-3-3-3-jcids-2025}.
  \item[\textbf{Personnel.}] People required to accomplish the mission~\autocite{edo-3-3-3-jcids-2025}.
  \item[\textbf{Facilities.}] Property, installations, and industrial facilities supporting forces~\autocite{edo-3-3-3-jcids-2025}.
  \item[\textbf{Policy.}] Policies that impact military operations~\autocite{edo-3-3-3-jcids-2025}.
\end{description}

\subsection{Capability Documents}
\begin{description}
  \item[\textbf{Purpose.}] Capability documents communicate operator needs to acquisition, test, and resourcing communities~\autocite{edo-3-3-3-jcids-2025}.
  \item[\textbf{Primary documents.}] \ac{icd} (approved prior to \ac{mdd}) and \ac{cdd} (approved prior to Milestone~B); \ac{cdd} updates support production~\autocite{edo-3-3-3-jcids-2025}.
  \item[\textbf{AoA support.}] \ac{aoa} is performed during the \ac{msa} phase as part of the requirements-to-acquisition handoff~\autocite{edo-3-3-3-jcids-2025}.
\end{description}

\subsection{Initial Capabilities Document (ICD)}
\begin{description}
  \item[\textbf{Purpose.}] Summarizes the \ac{cba} and justifies requirements to resolve capability gaps within a timeframe, including possible solutions for \ac{msa} and \ac{tmrr}~\autocite{edo-3-3-3-jcids-2025}.
  \item[\textbf{Inputs.}] Summarizes \ac{dotmlpfp} analysis and capability gaps~\autocite{edo-3-3-3-jcids-2025}.
  \item[\textbf{Supports.}] \ac{mdd}, \ac{aoa}, Acquisition Strategy, and Milestone~A~\autocite{edo-3-3-3-jcids-2025}.
\end{description}

\subsection{Capability Development Document (CDD)}
\begin{description}
  \item[\textbf{Purpose.}] Defines authoritative, measurable, and testable capabilities for \ac{emd} and the program technical baseline~\autocite{edo-3-3-3-jcids-2025}.
  \item[\textbf{Affordability.}] Captures information needed to deliver an affordable and supportable capability across increments~\autocite{edo-3-3-3-jcids-2025}.
  \item[\textbf{Performance.}] Specifies performance attributes, including \acp{kpp}, needed to design systems that close the gaps in the \ac{icd}~\autocite{edo-3-3-3-jcids-2025}.
  \item[\textbf{Updates.}] Can be updated near end of development to refine production attributes, quantities, and subsequent increments~\autocite{edo-3-3-3-jcids-2025}.
\end{description}

\subsection{CDD Update and Production Requirements}
\begin{description}
  \item[\textbf{Policy change.}] The 2018 \ac{jcids} manual removed the requirement for a separate \ac{cpd}; a \ac{cdd} update supports production needs~\autocite{edo-3-3-3-jcids-2025}.
  \item[\textbf{Legacy CPD intent.}] Defined capabilities for Production and Deployment, refined \ac{cdd} thresholds/objectives, and was approved before Milestone~C after \ac{cdr}~\autocite{edo-3-3-3-jcids-2025}.
  \item[\textbf{Legacy production role.}] A legacy \ac{cpd} defined the production baseline requirements used for test and evaluation, establishing the threshold performance attributes that had to be demonstrated before Milestone~C~\autocite{DoDI5000-85,edo-3-3-3-jcids-2025}.
  \item[\textbf{CDD update function.}] Modern \ac{cdd} updates incorporate production attributes, quantities, and supportability considerations that traditionally appeared in separate \ac{cpd} documents~\autocite{edo-3-3-3-jcids-2025}.
\end{description}

\subsection{JCIDS Within the Acquisition Framework}
\begin{figure}[htb]
  \centering
  \includegraphics[width=0.9\linewidth]{JCIDSinDAS.png}
  \caption[JCIDS Within the DAS]{Placement of \ac{jcids} documents within \ac{mca} phases. \srcCite{edo-3-3-3-jcids-2025}.}
  \label{fig:jcids-within-das}
\end{figure}

\begin{description}
  \item[\textbf{ICD.}] Defines capability gaps and supports \ac{mdd}, \ac{msa}, and \ac{tmrr} (no \acp{kpp})~\autocite{edo-3-3-3-jcids-2025}.
  \item[\textbf{CDD.}] Introduces \acp{kpp}, provides cost estimates, and guides \ac{emd} with measurable, testable capability definitions~\autocite{edo-3-3-3-jcids-2025}.
  \item[\textbf{CDD update.}] Refines \acp{kpp} and performance attributes, updates cost estimates, and guides Production and Deployment~\autocite{edo-3-3-3-jcids-2025}.
  \item[\textbf{Team effort.}] Requirements development must involve the warfighter and all stakeholders to produce executable technical requirements~\autocite{edo-3-3-3-jcids-2025}.
\end{description}

\subsection{Key Players}
\begin{description}
  \item[\textbf{Gatekeeper (J-8/DDR).}] Manages document flow, assigns \ac{jsd}, and designates lead/supporting \acp{fcb} for staffing and validation~\autocite{edo-3-3-3-jcids-2025}.
  \item[\textbf{Chairman of the Joint Chiefs of Staff.}] The \ac{cjcs} is the principal military adviser to the President, \ac{secdef}, and the \ac{nsc}; sets the agenda and presides over the Joint Chiefs of Staff~\autocite{edo-3-3-3-jcids-2025}.
  \item[\textbf{Joint Requirements Oversight Council.}] The \ac{jroc} is the process owner; validates requirements, prioritizes joint capabilities, and validates \acp{kpp} unless delegated~\autocite{edo-3-3-3-jcids-2025}.
  \item[\textbf{Joint Capabilities Board.}] The \ac{jcb} supports the \ac{jroc}, reviews documents, adjudicates issues, and validates \acp{kpp} for \ac{jcb}-interest documents~\autocite{edo-3-3-3-jcids-2025}.
  \item[\textbf{Functional Capabilities Boards.}] The \acp{fcb} are permanent bodies that analyze, prioritize, and integrate stakeholder views across functional areas; include \ac{ccmd} participation~\autocite{edo-3-3-3-jcids-2025}.
\end{description}

\subsection{Joint Staff J-6 (C4/Cyber)}
\begin{description}
  \item[\textbf{Net-ready validation.}] Validates the Net-Ready performance attribute (\ac{kpp} or \ac{ksa}) for \ac{jroc}/\ac{jcb}-interest documents~\autocite{edo-3-3-3-jcids-2025}.
  \item[\textbf{Interoperability certification.}] Performs IT and national security system interoperability and supportability certifications; \ac{jitc} issues certification before full-rate production~\autocite{edo-3-3-3-jcids-2025}.
\end{description}

\subsection{Key Performance Parameters and Performance Attributes}
\begin{description}
  \item[\textbf{Definition.}] \acp{kpp} and related performance attributes are system attributes critical to delivering effective capability and are documented in the \ac{cdd} and \ac{apb}~\autocite{DoDI5000-85,edo-3-3-2-acq-milestones-phases-2025}.
  \item[\textbf{Current governance.}] Under the current JFRP, the \ac{jroc} remains the senior joint requirements body, but Joint Staff review is focused on Joint Force Design, Joint Capability Integration, Combatant Command requirements, and by-exception reviews; Service and Component requirements are not automatically held for legacy \ac{jcids} validation before acquisition execution~\autocite{CJCSM5123-01}.
  \item[\textbf{Traceability.}] Programs still need clean traceability from validated capability needs to \acp{kpp}, \acp{ksa}, other performance attributes, verification methods, and decision evidence~\autocite{DoDI5000-85,edo-3-3-3-jcids-2025}.
\end{description}

\subsection{KPP vs KSA vs APA Comparison}
\begin{longtblr}[
    label   = {tab:kpp-ksa-apa-comparison},
    caption = {KPP vs KSA vs APA Comparison},
    entry   = {KPP vs KSA vs APA Comparison},
    remark{\tSource} = {\tabCite{CJCSM5123-01,DoDI5000-85,edo-3-3-3-jcids-2025}}
  ]
  {colspec = {@{} X[1.2] X[2.0] X[1.8] X[2.2] X[1.8] @{}}}
  \toprule
  {Parameter Type} & {Definition} & {Governance Level} & {Impact if Unmet} & {Navy Example} \\
  \midrule
  \textbf{Key Performance Parameter (KPP)} & Critical capability attribute tied to the validated capability need and acquisition baseline & JROC/JFRP review when tied to joint-force requirements or by-exception joint impact; Component validation for Service-specific programs & Program stoppage; statutory reporting required; Milestone decisions blocked & Illustrative Aegis Weapon System example: defended-area detection coverage sufficient to support raid engagement timelines \\
  \midrule
  \textbf{Key System Attribute (KSA)} & Important performance attribute that supports KPP achievement and operational suitability & Joint Staff or functional review where the attribute has joint/interoperability impact; Component validation for program-specific attributes & Interoperability certification delays; reduced operational effectiveness; potential re-baseline & Illustrative Aegis Weapon System example: identification data exchange responsive enough to support engagement coordination \\
  \midrule
  \textbf{Additional Performance Attribute (APA)} & Performance attribute required for full system capability but not critical enough for KPP & Component validation; documented in CDD/APB & Reduced test priority; capability gaps if not met; schedule impacts & Illustrative Aegis Weapon System example: electronic-protection behavior that preserves mission effectiveness in representative threat environments \\
  \bottomrule
\end{longtblr}

\subsection{IT Acquisition Requirements}
\begin{description}
  \item[\textbf{Net-Ready attribute.}] Documented as a \ac{kpp} or \ac{ksa} to ensure IT supports operations, is manageable on networks, and exchanges information effectively~\autocite{edo-3-3-3-jcids-2025}.
  \item[\textbf{Architecture compliance.}] Systems must align to \ac{dodaf}; architecture views are mandatory parts of the \ac{cdd} and must be consistent with \ac{disr} standards~\autocite{edo-3-3-3-jcids-2025}.
\end{description}

\subsection{Importance of Interoperability}
\begin{description}
  \item[\textbf{Definition.}] Ability of systems, units, or forces to provide and accept services and operate effectively with U.S. and coalition partners~\autocite{edo-3-3-3-jcids-2025}.
  \item[\textbf{Operational risk.}] Poor interoperability increases risk to people and equipment and reduces operational effectiveness~\autocite{edo-3-3-3-jcids-2025}.
  \item[\textbf{Cost discipline.}] Stovepiped systems are costly and hinder joint operations; interoperability is a core \ac{jcids} validation priority~\autocite{edo-3-3-3-jcids-2025}.
\end{description}


\subsection{Validation Authority}
\begin{description}
  \item[\textbf{Joint Requirements Oversight Council.}] Under the JFRP, the \ac{jroc} remains the senior joint requirements body and reviews prioritized Joint Force and Combatant Command requirement issues, but it is not the default validator for every Service-generated acquisition requirement document~\autocite{CJCSM5123-01}.
  \item[\textbf{Joint Capabilities Board.}] The \ac{jcb} supports \ac{jroc} deliberations and adjudicates staffing issues for Joint Capability Integration and by-exception reviews directed by the Joint Staff~\autocite{CJCSM5123-01}.
  \item[\textbf{Component sponsor.}] Services and Components generate and validate their own requirements; Joint Capability Integration occurs after Component validation and in parallel with acquisition execution when joint impact warrants it~\autocite{CJCSM5123-01}.
\end{description}

\subsection{Document Review and Staffing}
\begin{description}
  \item[\textbf{Submission.}] Under legacy JCIDS teaching flow, sponsors submit documents to Joint Staff J-8 for staffing; under current JFRP, this routing is selective and focused on joint-force integration priorities rather than universal staffing of Service documents~\autocite{CJCSM5123-01,edo-3-3-3-jcids-2025}.
  \item[\textbf{JSD assignment.}] J-8 assigns staffing and review mechanisms according to joint impact and JFRP governance decisions~\autocite{CJCSM5123-01}.
  \item[\textbf{FCB review.}] Lead \ac{fcb} analysis still informs joint integration decisions when a requirement has cross-functional or joint implications~\autocite{CJCSM5123-01}.
  \item[\textbf{Validation.}] Board answer: state Component validation first, then explain Joint Staff/JROC integration review as required by scope and joint impact~\autocite{CJCSM5123-01}.
\end{description}

\subsection{Quick Review}
\begin{itemize}
  \item What system institutionalizes a joint capability need and sponsors further evaluation of a materiel solution~\autocite{edo-3-3-3-jcids-2025}? \textbf{Answer:} \ac{jcids}.
  \item Which acquisition phases use the \ac{icd} and \ac{cdd}~\autocite{edo-3-3-3-jcids-2025}? \textbf{Answer:} \ac{msa}/\ac{tmrr}, \ac{emd}, and Production and Deployment.
  \item Which document defines capabilities and is transformed into a program start~\autocite{edo-3-3-3-jcids-2025}? \textbf{Answer:} \ac{cdd}, with Milestone~B approval by the \ac{mdauthority}.
  \item JROC and the Joint Staff ensure new capabilities are conceived and developed in what context~\autocite{edo-3-3-3-jcids-2025}? \textbf{Answer:} Joint warfighting context.
  \item Warfighter gap identification is driven by Combatant Commander participation in what body~\autocite{edo-3-3-3-jcids-2025}? \textbf{Answer:} \acp{fcb}.
\end{itemize}

%====================
% End of file
%====================
\IfSubfilesClassLoaded{
  \RenewDocumentCommand{\entryname}{}{\textbf{\color{Modern} Acronym}}
  \RenewDocumentCommand{\descriptionname}{}{\textbf{\color{Modern} Definition}}
  \printnoidxglossary[
    type=\acronymtype,
    title=Acronyms,
    style=long-booktabs
  ]}{}
\subsection{Coursebook cue: legacy JCIDS language}

Coursebook 3.3.3 uses JCIDS-era vocabulary to teach capability gaps,
requirements documents, and validation.  For a board answer, state the
coursebook concept first, then add the current-policy caveat when relevant:
legacy JCIDS terms remain useful for explaining need, gap, validation, and
traceability, but current Joint Staff requirements reform may change document
names, routing, and review mechanics
\autocite{edo-3-3-3-jcids-2025}.

\end{document}
